\documentclass[../../main.tex]{subfiles}
\begin{document}

\paragraph{【乙】}
$xyz$空間で$O(0,0,0)$，$A(3,0,0)$，$B(3,2,0)$，$C(0,2,0)$，
$D(0,0,4)$，$E(3,0,4)$，$F(3,2,4)$，$G(0,2,4)$を頂点とする
直方体$OABC-DEFG$を考える．
辺$AE$を$s:1-s$に内分する点を$P$，
辺$CG$を$t:1-t$に内分する点を$Q$とおく．
ただし$0<s<1$，$0<t<1$とする．
$D$を通り，$O$，$P$，$Q$を含む平面に垂直な直線が線分$AC$(両端を含む)と交わるような
$s$，$t$のみたす条件を求めよ．

\paragraph{【甲】}
次の各問にそれぞれ答えよ．
\begin{enumerate}
  \item 正の数$a$に対して$xyz$空間で$O(0,0,0)$，$A(3,0,0)$，$B(3,2,0)$，$C(0,2,0)$，$D(0,0,a)$，$E(3,0,a)$，$F(3,2,a)$，$G(0,2,a)$を頂点とする直方体$OABC-DEFG$を考える．$D$を通り，3つの頂点$O$，$E$，$G$を含む平面に垂直な直線が辺$BC$（両端を含む）と点$P$で交わるとき，$a$の値と$P$の座標を求めよ．
  \item 白球と赤球の入った袋から2個の球を同時に取り出すゲームを考える．取り出した2球がともに白球ならば「成功」でゲームを終了し，そうでないときは「失敗」とし，取り出した2球に赤球を1個加えた3個の球を袋にもどしてゲームを続けるものとする．最初に白球が2個，赤球が1個袋に入っていたとき，$n-1$回まで失敗し$n$回目に成功する確率を求めよ．ただし$n \geqq 2$とする．
\end{enumerate}

\end{document}
